\section*{अध्याय \ref{ch:g10:storms-and-cyclones} --- तूफ़ान और उष्णकटिबंधीय चक्रवात}

\begin{solution}{exo:g10:storms-and-cyclones:1}
\omterm{def:g10:storms-and-cyclones:thunderstorm}{तड़ित-झंझा} किसी कपासी वर्षी मेघ से उत्पन्न ऐसा तूफ़ान है जिसमें \omterm{def:g10:storms-and-cyclones:lightning}{तड़ित} और
गरज होती है, सामान्यतः तेज़ बौछारें और झोंके, कभी-कभी ओले। \omterm{def:g10:storms-and-cyclones:lightning}{तड़ित} एक विशाल
विद्युत चिंगारी है जो तूफ़ान के आवेशित क्षेत्रों को अनावेशित करती है। गरज
\omterm{def:g10:storms-and-cyclones:lightning}{तड़ित-नली} की उस हवा की ध्वनि है जो इतनी अचानक गरम होती है कि एक आघात के साथ
फैल जाती है।
\end{solution}

\begin{solution}{exo:g10:storms-and-cyclones:2}
\omterm{def:g10:storms-and-cyclones:tropical-cyclone}{उष्णकटिबंधीय चक्रवात} सैकड़ों किलोमीटर चौड़ा ऐसा तूफ़ान है जो गरम उष्णकटिबंधीय
समुद्र के ऊपर निम्न दाब के एक केंद्र के चारों ओर बनता है, जिसमें पवनें सर्पिल रूप में भीतर घूमती हैं
और \omterm{def:g10:storms-and-cyclones:thunderstorm}{तड़ित-झंझाओं} की पट्टियाँ होती हैं। \omterm{def:g10:storms-and-cyclones:tropical-cyclone}{आँख} उसका शांत, लगभग मेघरहित केंद्र है;
\omterm{def:g10:storms-and-cyclones:tropical-cyclone}{आँख की दीवार} \omterm{def:g10:storms-and-cyclones:tropical-cyclone}{आँख} के चारों ओर ऊँची \omterm{def:g10:storms-and-cyclones:thunderstorm}{तड़ित-झंझाओं} का वलय है। पवनें
\omterm{def:g10:storms-and-cyclones:tropical-cyclone}{आँख की दीवार} में सबसे तेज़ होती हैं।
\end{solution}

\begin{solution}{exo:g10:storms-and-cyclones:3}
\num{70}~नॉट: श्रेणी 1। \num{120}~नॉट: श्रेणी 4।
\num{140}~नॉट: श्रेणी 5। अंतिम दो प्रमुख हरिकेन हैं (श्रेणी 3
या अधिक)।
\end{solution}

\begin{solution}{exo:g10:storms-and-cyclones:4}
$343 \times 9 \approx \qty{3.1}{km}$, फिर $343 \times 15 \approx
\qty{5.1}{km}$। दूसरी चमक अधिक दूर है: तूफ़ान दूर जा रहा है।
\end{solution}

\begin{solution}{exo:g10:storms-and-cyclones:5}
\omterm{def:g10:storms-and-cyclones:lightning}{तड़ित} आसपास के सबसे ऊँचे बिंदु पर गिरने की प्रवृत्ति रखती है: खेत में अकेला खड़ा पेड़
वही बिंदु है, और प्रहार की धारा उससे होकर नीचे उतरती है और उसके चारों ओर की ज़मीन में
फैल जाती है। किसी इमारत में या किसी बंद कार में जाइए, और अंतिम गरज के काफ़ी देर
बाद तक वहीं रहिए।
\end{solution}

\begin{solution}{exo:g10:storms-and-cyclones:6}
एक नॉट $1852/3600 \approx \qty{0.514}{m/s}$ है। \num{64}~नॉट: $64 \times
0.514 \approx \qty{32.9}{m/s}$, और $32.9 \times 3.6 \approx
\qty{119}{km/h}$। \num{137}~नॉट: $\qty{70.5}{m/s}$, लगभग
\qty{254}{km/h} (आधिकारिक तालिका, किसी दूसरी इकाई से पूर्णांकित करके,
\qty{252}{km/h} देती है)।
\end{solution}

\begin{solution}{exo:g10:storms-and-cyclones:7}
$44 \times 86\,400 \approx \num{3.8e6}$ चमकें प्रतिदिन, और $3.8\times10^6
\times 365.25 \approx \num{1.4e9}$, लगभग 1.4 अरब प्रति वर्ष।
\end{solution}

\begin{solution}{exo:g10:storms-and-cyclones:8}
$\Delta p = 1010 - 935 = \qty{75}{hPa}$: लगभग \qty{75}{cm} का उभार।
\end{solution}

\begin{solution}{exo:g10:storms-and-cyclones:9}
वे गरम उष्णकटिबंधीय महासागरों के ऊपर बनते हैं, अधिकतर दोनों गोलार्धों में \qty{10}{\degree} और
\qty{30}{\degree} अक्षांश के बीच, और दक्षिणी अटलांटिक में बहुत कम।
वे पश्चिम की ओर बहते हैं, फिर मुड़कर ध्रुवों की ओर चले जाते हैं, जहाँ वे ठंडे जल के ऊपर
कमज़ोर हो जाते हैं। वे भूमध्य रेखा पर इसलिए नहीं होते कि वहाँ
पृथ्वी का घूर्णन भीतर बहती हवा को कोई मोड़ नहीं देता: उसके बिना हवा सीधे
निम्न दाब में बहकर उसे भर देती है, और कोई घूमता तूफ़ान शुरू नहीं हो सकता।
\end{solution}

\begin{solution}{exo:g10:storms-and-cyclones:10}
$1031 - 1.56 \times 70 \approx \qty{922}{hPa}$; $1031 - 1.56 \times 40
\approx \qty{969}{hPa}$।
\end{solution}

\begin{solution}{exo:g10:storms-and-cyclones:11}
स्थल के ऊपर तूफ़ान वाष्पित होते जल की अपनी आपूर्ति खो देता है, और उसके साथ वह
\omterm{def:g10:earth-energy-budget:latent-heat}{गुप्त ऊष्मा} भी जो उसे चलाती है; खुरदरी सतह भी पवनों को धीमा कर देती है। ठंडे
जल के ऊपर वाष्पन बहुत कमज़ोर होता है, क्योंकि ठंडी हवा संतृप्ति पर कहीं कम
जलवाष्प धारण करती है: \omterm{def:g10:earth-energy-budget:latent-heat}{गुप्त ऊष्मा} की आपूर्ति घटती जाती है।
\end{solution}

\begin{solution}{exo:g10:storms-and-cyclones:12}
\qty{0}{\degreeCelsius} पर: $\sqrt{1.4 \times 8.314 \times 273.15/0.02897}
\approx \qty{331}{m/s}$; \qty{30}{\degreeCelsius} पर, $T =
\qty{303.15}{K}$ के साथ: \qty{349}{m/s}। \qty{9}{s} के लिए: $9 \times 331 \approx
\qty{2980}{m}$ और $9 \times 349 \approx \qty{3140}{m}$, लगभग
\qty{160}{m} का अंतर, लगभग \qty{5}{\percent}। नियम \qty{3}{km} देता है:
यह जानने के लिए पर्याप्त कि तूफ़ान कितना पास है, और इसका काम इतना ही है।
\end{solution}

\begin{solution}{exo:g10:storms-and-cyclones:13}
$\Delta p = 1010 - 940 = \qty{70}{hPa}$: उभार लगभग \qty{0.7}{m} है,
पवन का भाग $4.0 - 0.7 = \qty{3.3}{m}$, चार गुना से भी अधिक। ऊँचे
ज्वार पर जल माध्य समुद्र तल से $4.0 + 1.5 = \qty{5.5}{m}$ ऊपर तक पहुँचता है।
\end{solution}

\begin{solution}{exo:g10:storms-and-cyclones:14}
\omterm{def:g10:storms-and-cyclones:tornado}{टॉरनेडो}: एक चक्कर $2\pi \times 100 \approx \qty{628}{m}$ है, जो
$628/90 \approx \qty{7}{s}$ में तय होता है। हरिकेन: $2\pi \times 30\,000 \approx
\qty{188}{km}$, जो $188\,000/60 \approx \qty{3140}{s}$ में, लगभग
\qty{52}{min} में, तय होता है। \omterm{def:g10:storms-and-cyclones:tornado}{टॉरनेडो} लगभग 450 गुना तेज़ घूमता है: वह एक छोटा, बहुत
तेज़ भँवर है, हरिकेन एक विशाल और धीमा भँवर।
\end{solution}

\begin{solution}{exo:g10:storms-and-cyclones:15}
$24/3888 \approx \qty{0.6}{\percent}$; $414/3888 \approx
\qty{10.6}{\percent}$; $121/3888 \approx \qty{3.1}{\percent}$। यदि अक्षांश का
कोई महत्त्व न होता, तो लगभग \qty{9}{\percent} तूफ़ान, कोई 350, भूमध्य रेखा के
\qty{5}{\degree} के भीतर शुरू होते; केवल 24 होते हैं, लगभग पंद्रह गुना
कम। भूमध्य रेखा के पास पृथ्वी का घूर्णन घुमाव शुरू करने के लिए बहुत कमज़ोर है।
\end{solution}

\begin{solution}{pb:g10:storms-and-cyclones:1}
\textbf{1.} $\pi \times (665\,000)^2 \approx \qty{1.39e12}{m^2}$।

\textbf{2.} $0.015 \times 1.39\times10^{12} \approx \qty{2.08e10}{m^3}$।

\textbf{3.} एक घन किलोमीटर \qty{e9}{m^3} है: लगभग \qty{21}{km^3}।

\textbf{4.} $2.08\times10^{10} \times 1000 \approx \qty{2.08e13}{kg}$।

\textbf{5.} $2.08\times10^{13} \times 2.5\times10^6 \approx \qty{5.2e19}{J}$।

\textbf{6.} $5.2\times10^{19}/86\,400 \approx \qty{6.0e14}{W}$।

\textbf{7.} $6.0\times10^{14}/1.22\times10^{17} \approx 0.005$, लगभग
\qty{0.5}{\percent}।

\textbf{8.} $240 \times 1.39\times10^{12} \approx \qty{3.3e14}{W}$। तूफ़ान
अपने पूरे क्षेत्रफल द्वारा अवशोषित सूर्य-प्रकाश की लगभग दुगुनी ऊष्मा मुक्त करता है।

\textbf{9.} तूफ़ान अपने ऊपर पड़ने वाले सूर्य-प्रकाश पर नहीं चलता: वह जलवाष्प के रूप
में संचित ऊष्मा मुक्त करता है, जो गरम समुद्र के कहीं अधिक विस्तृत क्षेत्र से वाष्पित हुई
और पवनों द्वारा भीतर खींची गई। तूफ़ान कहीं अधिक बड़े क्षेत्र में एकत्र
ऊर्जा को सांद्रित करता है।

\textbf{10.} वर्षा जितना ही द्रव्यमान, लगभग \qty{2.1e13}{kg} प्रतिदिन (भीतर लाई गई
वाष्प का नवीकरण उतनी ही तेज़ी से होना चाहिए जितनी तेज़ी से वह बरसती है)।

\textbf{11.} $6.0\times10^{14}/82 \approx \qty{7.3e12}{m^2}$, तूफ़ान के क्षेत्रफल का लगभग
पाँच गुना: तूफ़ान को कहीं अधिक विस्तृत क्षेत्र से वाष्प खींचनी पड़ती है, और
उसकी पवनें अपने नीचे के समुद्र को औसत से कहीं अधिक तेज़ी से वाष्पित कराती हैं।

\textbf{12.} \qty{28}{\degreeCelsius} पर \qty{24.1}{g/kg}, जबकि
\qty{20}{\degreeCelsius} पर \qty{14.7}{g/kg}: $24.1/14.7 \approx 1.64$, लगभग
\qty{64}{\percent} अधिक।

\textbf{13.} गरम जल तेज़ी से वाष्पित होता है और गरम हवा कहीं अधिक वाष्प धारण करती है:
भीतर बहती हवा मुक्त करने के लिए कहीं अधिक \omterm{def:g10:earth-energy-budget:latent-heat}{गुप्त ऊष्मा} लाती है, और संतृप्त होकर
वह बेरोक-टोक ऊपर उठती है। ठंडे समुद्रों के ऊपर आपूर्ति इंजन को चलाने के लिए बहुत
कम है।

\textbf{14.} स्थल के ऊपर वाष्पित होने के लिए कोई समुद्र नहीं होता: \omterm{def:g10:earth-energy-budget:latent-heat}{गुप्त ऊष्मा} की
आपूर्ति रुक जाती है, और खुरदरी ज़मीन पवनों को धीमा कर देती है।

\textbf{15.} $1.5\times10^{12} \times 86\,400 \approx \qty{1.3e17}{J}$।

\textbf{16.} $96 \times 1852/3600 \approx \qty{49.4}{m/s}$।

\textbf{17.} $\frac12 \times 1.16 \times 49.4^2 \approx \qty{1.4e3}{J}$।

\textbf{18.} $5.2\times10^{19}/1.3\times10^{17} = 400$: \omterm{def:g10:earth-energy-budget:latent-heat}{गुप्त ऊष्मा}
पवनों को दी गई ऊर्जा की लगभग 400 गुना है।

\textbf{19.} इसका अधिकांश भाग तूफ़ान की ऊपरी हवा को गरम करता है, जो शीर्ष पर फैल
जाती है और ऊष्मा को अवरक्त विकिरण के रूप में अंतरिक्ष को दे देती है; पवनें स्वयं
समुद्र और ज़मीन से घर्षण में अपनी ऊर्जा खो देती हैं।

\textbf{20.} $1/400 = \qty{0.25}{\percent}$ (स्क्रिप्ट
402 का अनुपात देती है): \textbf{संघनन से मुक्त हुई ऊष्मा का लगभग 400 में एक भाग
पवन बनता है}।
\end{solution}
